\scanpage{149}
\margin{148}
\mspage{9.}

\uline{Théorème 1.}- \uline{Soit} $X$ \uline{un préschéma. Considérons les sous-catégories pleines} $E$ \uline{de} $\overset{\inserted{\wedge}}{(\mathrm{Sch})_{/X}}$ \uline{qui sont des} $\underline{U}$-\uline{topos et dont les objets} $F$ \uline{satisfont aux axiomes}

(A) \uline{Pour tout} $X$-\uline{morphisme} $S\inserted{\xrightarrow{\varphi}}S'$ \uline{entre deux} $X$-\uline{préschémas affines dans le sens absolu tel que} $\inserted{\varphi}$ \uline{induise un homéomorphisme sur les espaces topologiques sous-jacentes, le morphisme évident}
\diagram{\begin{tikzcd}[ampersand replacement=\&]
F(S')\arrow[r,"F(\inserted{\varphi})"]\&F(S)
\end{tikzcd}}
\uline{est un isomorphisme,}

(B) \uline{Pour tout} $\underline{U}$-\uline{système projectif filtrant de préschémas} $\inserted{\{}S_{\inserted{\alpha}}\inserted{\}}$ \uline{où les} $S$ \uline{sont affines au-dessus d'un ouvert affine de} $X$ \uline{l'application évidente}
\diagram{\begin{tikzcd}[ampersand replacement=\&]
F(\varprojlim S_\alpha)\&\varinjlim F(S_\alpha)\arrow[l]
\end{tikzcd}}
\uline{est un isomorphisme,}

\uline{et restreignons-nous aux} $E$ \uline{qui de plus contiennent les objets représentables satisfaisant aux conditions} (A) \uline{et} (B). \uline{Parmi ces topos il y a un qui est} minimum \uline{i.e.\ est contenu dans tout autre} $E$. \uline{Ce topos} (évidemment unique ) \uline{est équivalent à} $\overset{\inserted{\sim}}{\operatorname{Ouv}(X)}$.

La démonstration du théorème 1 s'appuiera sur le lemme suivant qui sera démontré dans le \S\ 3.

\uline{Lemme 2.}- \uline{La condition nécessaire et suffisante pour qu'un objet} $Y/X\inserted{\in}(\mathrm{Sch})_{/X}\inserted{\hookrightarrow}(\mathrm{Sch})_{/X}$ \uline{satisfasse aux conditions} (A) \uline{et} (B) \uline{du théorème} 1, \uline{c'est que} $Y\inserted{\xrightarrow{f}}X$ \uline{soit un isomorphisme local.}
(Ceci signifie que chaque point de $Y$ admet un voisinage ouvert $U$ tel que $f\inserted{|}U:U\inserted{\longrightarrow}X$ soit une immersion ouverte \inserted{[}EGA I.4.5\inserted{]} )

\uline{Démonstration du théorème 1.}- Soit $\operatorname{Isloc}(X)$ la catégorie des isomorphismes locaux $Y/X$. Pour démontrer le théorème 1, il suffit de voir qu'il y a une équivalence de catégories
\diagram{\begin{tikzcd}[ampersand replacement=\&]
\overset{\inserted{\sim}}{\operatorname{Ouv}(X)}\arrow[r]\&\operatorname{Isloc}(X).
\end{tikzcd}}
